\documentclass[11pt]{article}
\usepackage{ochamath}
\subjectcolor{2F5DA8}
\setsheet{線形代数}{対称行列と直交対角化　演習}{https://math.ochanote.com/linear-algebra/symmetric-matrices/}
\begin{document}
\makesheethead{解答}

\problem[直交対角化]
\begin{ans}
固有値6と2に対応する正規化ベクトルは $(1,1)^{\mathsf T}/\sqrt2$、$(1,-1)^{\mathsf T}/\sqrt2$。\[Q=\frac1{\sqrt2}\begin{pmatrix}1&1\\1&-1\end{pmatrix},\quad D=\operatorname{diag}(6,2).\] $Q^{\mathsf T}Q=I,\ AQ=QD$。
\end{ans}
\point{固有ベクトルは長さ1へ正規化する。}

\problem[重複固有値]
\begin{ans}
内積は1で直交しない。どちらも固有値1だが、定理は正規直交基底を選べることを述べる。例えば標準基底は直交するため矛盾しない。
\end{ans}
\point{異なる固有値の直交性と、同じ固有空間での基底選択を区別する。}
\newpage

\problem[長さの保存]
\begin{ans}
座標は $(\sqrt2,\sqrt2)^{\mathsf T}$。2乗の和は $2+2=4$ なので長さは2で、元の $(2,0)^{\mathsf T}$ の長さと同じ。
\end{ans}
\point{直交基底の座標では長さの計算も同じ形式になる。}

\problem[剛性のモード]
\begin{ans}
$(1,1)^{\mathsf T}$ は固有値1、$(1,-1)^{\mathsf T}$ は固有値3。$\omega^2$ が固有値なので、角振動数は1と $\sqrt3$。逆方向の方が高い周波数となる。
\end{ans}
\point{剛性の固有値と周波数そのものを取り違えない。}
\end{document}
